\documentclass[10pt]{article}

\usepackage[margin=1in]{geometry}
\usepackage{amsmath, amssymb, amsthm}
\usepackage{algorithm, algpseudocode}
\usepackage{booktabs}
\usepackage{hyperref}
\usepackage{microtype}

\title{Task-Global Variance Normalization for GRPO:\\
Reallocating Gradient Budget by Within-Group Reward Dispersion}

\author{Anonymous}
\date{}

\begin{document}
\maketitle

\begin{abstract}
Group Relative Policy Optimization (GRPO) normalizes per-response rewards by the
empirical standard deviation $\sigma_G$ computed within the response group sampled
for a single prompt. This forces every group's advantages to have unit variance
regardless of how informative the group is, with two symmetric pathologies: when
$\sigma_G$ is small (responses cluster tightly), sub-reward-model-noise differences
are amplified into apparently-strong gradients; when $\sigma_G$ is large (responses
disagree meaningfully), genuine signal is compressed to the same magnitude as the
noise. We propose replacing the per-group denominator $\sigma_G$ with a
task-conditional, EMA-smoothed global estimate $\hat{\sigma}_{\text{global}}(T)$,
followed by a hard clip on the resulting advantage. The change is one line of code,
adds negligible compute, and makes the per-group gradient contribution proportional
to $\sigma_G / \hat{\sigma}_{\text{global}}(T)$---automatically reallocating
gradient budget toward groups whose response dispersion carries learning signal.
On the MATH benchmark we observe a difficulty-monotonic effect consistent with the
mechanism: accuracy decreases on saturated easy problems and increases on the
hardest tier.
\end{abstract}

\section{The Group Variance Pathology}

\paragraph{Setup.}
GRPO~\citep{shao2024deepseekmath} samples a group $G = \{o_1,\dots,o_n\}$ of
responses for a single prompt $q$, scores each with a reward model to obtain
$R_i$, and computes the advantage
\begin{equation}
A_i \;=\; \frac{R_i - \mu_G}{\sigma_G + \epsilon},
\qquad
\mu_G = \tfrac{1}{n}\sum_j R_j,
\quad
\sigma_G = \sqrt{\tfrac{1}{n}\sum_j (R_j - \mu_G)^2}.
\label{eq:grpo}
\end{equation}
The advantage $A_i$ is a scalar, but in the policy-gradient update
$\nabla_\theta\mathcal{L} = -\sum_i A_i\,\nabla_\theta\log\pi_\theta(o_i\!\mid\!q)$
it serves as the per-sample weight on a high-dimensional gradient vector. Its
\emph{magnitude} therefore directly controls the effective learning step taken on
behalf of each response.

\paragraph{The pathology, two symmetric faces.}
Equation~\eqref{eq:grpo} forces $A_i$ to have unit standard deviation \emph{within
every group}, regardless of how informative that group is. This is a problem in
both directions.

\medskip
\noindent\textbf{Case A: low-$\sigma_G$ groups (mastered or low-resolution tasks).}
Consider a creative-writing prompt where all eight rubric-judged completions are
reasonable, with rewards $\{7.2, 7.0, 7.1, 6.9, 7.3, 7.0, 7.1, 7.2\}$:
$\mu_G \approx 7.1,\ \sigma_G \approx 0.13$. The $0.2$-point gap between the
``best'' and the mean is well below the rubric judge's own grading-noise floor,
yet GRPO computes
\[
A_{\text{best}} = \frac{7.3 - 7.1}{0.13} \approx 1.54,
\]

\[
A_{\text{best}} = \frac{10 - 5.25}{3.77} \approx 1.26
\]

a unit-magnitude signal indistinguishable from the genuine $A=+1$ produced by a
correct-vs-wrong math group. The optimizer cannot tell the difference and updates
the policy as if the $7.3$-rated completion were meaningfully better than the
$7.0$-rated one. The same effect arises on prompts the policy has already
mastered---all sampled responses receive near-maximal reward, so $\sigma_G$
collapses and judge noise becomes the dominant gradient signal.

\medskip
\noindent\textbf{Case B: high-$\sigma_G$ groups (frontier-difficulty problems).}
Conversely, on a hard math prompt where four of eight sampled solutions are
correct ($R=100$) and four are wrong ($R=0$), $\mu_G=50$ and $\sigma_G=50$. GRPO
again normalizes to $A_i \in \{-1, +1\}$. But this group carries far more
learning signal than the writing group above---the policy has visible room to
improve, and the within-group disagreement is real. By compressing it to the
same unit magnitude, GRPO gives the difficult problem no more gradient weight
than the saturated easy one.

\paragraph{What the two cases share.}
GRPO spends the same gradient budget on every group regardless of how much
genuine information that group carries. Easy or mastered prompts receive too
much, hard or genuinely contested prompts receive too little, and the optimizer
has no signal to distinguish the two.

\section{Method: Task-Global Variance Normalization}

\paragraph{Core idea.}
Replace the group-local $\sigma_G$ with a task-conditional, slowly-updated global
estimate of reward dispersion:
\begin{equation}
A_i^{\text{TG}} \;=\; \frac{R_i - \mu_G}{\hat{\sigma}_{\text{global}}(T_i) + \epsilon}
\label{eq:tg}
\end{equation}
where $T_i \in \mathcal{T}$ is the task type of sample $i$ (math, writing, code,
\ldots) and $\hat{\sigma}_{\text{global}}(T)$ is maintained by an exponential
moving average over per-batch within-task variances:
\begin{equation}
\hat{\sigma}^2_{\text{global}}(T_i) \;\leftarrow\; \alpha\,\sigma^2_{\text{batch}}(T_i)
\;+\; (1-\alpha)\,\hat{\sigma}^2_{\text{global}}(T_i)
\qquad \alpha = 0.1.
\label{eq:ema}
\end{equation}
The numerator $R_i - \mu_G$ is unchanged; baseline subtraction still uses the
group mean. Only the scaling denominator is decoupled from the current group.

\paragraph{Why this fixes both cases simultaneously.}
Under~\eqref{eq:tg}, the magnitude of $A_i$ within a group is no longer
artificially rescaled to unit variance. Instead it satisfies
\begin{equation}
\|A^{\text{TG}}\|_{\text{group}}
\;\propto\;
\frac{\sigma_G}{\hat{\sigma}_{\text{global}}(T)}.
\label{eq:budget}
\end{equation}
A group's contribution to the gradient is now proportional to how dispersed
\emph{this} group is relative to the typical dispersion of its task. The
mechanism is symmetric:

\begin{itemize}\itemsep2pt
\item $\sigma_G \ll \hat{\sigma}_{\text{global}}(T)$: group is unusually
homogeneous (mastered prompt or low-resolution judge); advantages are
\emph{shrunk}, so the optimizer is not pushed to learn from sub-noise
differences.
\item $\sigma_G \approx \hat{\sigma}_{\text{global}}(T)$: typical group;
behaviour matches standard GRPO.
\item $\sigma_G \gg \hat{\sigma}_{\text{global}}(T)$: group is unusually
dispersed (frontier-difficulty prompt with real disagreement); advantages are
\emph{amplified}, concentrating learning effort where genuine signal exists.
\end{itemize}

For the writing group from Case~A, with $\hat{\sigma}_{\text{global}}(\text{Write})
\approx 1.5$ estimated over the writing task as a whole,
$A^{\text{TG}}_{\text{best}} = (7.3 - 7.1)/1.5 \approx 0.13$: the small reward gap
is correctly reflected as a small advantage. For a high-disagreement math group
with $\sigma_G = 50$ and $\hat{\sigma}_{\text{global}}(\text{Math}) \approx 30$,
$A^{\text{TG}} \approx \pm 1.67$ rather than $\pm 1.0$, giving the harder
prompt the larger gradient weight it warrants.

\paragraph{The new failure mode and its fix.}
Decoupling the denominator from the current group sacrifices one nice property of
standard GRPO: under the original normalization, an outlier reward (e.g.\ a
reward-model bug returning $R=5000$) inflates $\sigma_G$ and is therefore
self-absorbed. Under~\eqref{eq:tg}, $\hat{\sigma}_{\text{global}}$ has too much
inertia to react, and the spike passes through to the advantage. We close this
gap with a hard clip:
\begin{equation}
A_i^{\text{final}} \;=\; \mathrm{clip}\!\left(A_i^{\text{TG}},\,-C,\,C\right),
\qquad C = 10.
\label{eq:clip}
\end{equation}
Under a well-behaved reward model and Equation~\eqref{eq:tg},
$A^{\text{TG}}_i$ is approximately zero-mean with task-level standard deviation
$\approx 1$, so $|A^{\text{TG}}_i| > 10$ is a $10\sigma$ event---essentially
diagnostic of a malfunction rather than a real signal. Clipping at $C=10$
preserves the full ordinal ranking of legitimate responses while bounding the
worst-case per-sample contribution to the gradient.

\paragraph{Task identification.}
We assume access to a task tag $T_i$ either from the dataset or from a trivial
reward-magnitude rule when reward ranges across tasks are well-separated. The
method is agnostic to how $T_i$ is obtained.

\section{Empirical Signature on MATH}

We train a 7B-parameter Qwen-Instruct model with V-GRPO and compare against a
GRPO baseline on the MATH benchmark, broken down by difficulty level.
Table~\ref{tab:math} reports per-level accuracy.

\begin{table}[h]
\centering
\small
\begin{tabular}{lccr}
\toprule
\textbf{MATH Level} & \textbf{Baseline (\%)} & \textbf{V-GRPO (\%)} & $\Delta$ \\
\midrule
Level 1             & 80.37 & 80.64 & $+0.27$ \\
Level 2             & 68.78 & 68.37 & $-0.41$ \\
Level 3    & 47.21 & 48.26 & $+1.05$ \\
\bottomrule
\end{tabular}
\caption{V-GRPO vs.\ GRPO on the MATH benchmark, per difficulty level. The
pattern is monotonic in difficulty: V-GRPO loses on saturated easy problems
and gains on the hardest tier.}
\label{tab:math}
\end{table}

\paragraph{Reading the result through the mechanism.}
The pattern in Table~\ref{tab:math} is what
Equation~\eqref{eq:budget} predicts. Easy MATH problems are saturated for the
base policy: most sampled solutions are correct, $\sigma_G$ is small relative to
$\hat{\sigma}_{\text{global}}(\text{Math})$, and V-GRPO shrinks the advantages
that GRPO would have inflated. The policy stops over-rehearsing material it has
already mastered and accuracy drifts down on Level~1. Hard problems exhibit
genuine response disagreement, $\sigma_G$ approaches or exceeds
$\hat{\sigma}_{\text{global}}(\text{Math})$, and V-GRPO amplifies the
advantages---directing more learning effort to the regime where the policy can
still improve. Mid-difficulty levels, whose $\sigma_G$ tracks the task average,
move negligibly. The aggregate accuracy change is small ($-0.42$ averaged across
levels), but the difficulty-monotonic shape of the per-level deltas is the
direct empirical signature of gradient-budget reallocation.

\paragraph{What this is and is not evidence of.}
This single-task evaluation does not establish that V-GRPO improves overall
benchmark numbers. It does establish that the mechanism described in
Equation~\eqref{eq:budget} is operative: a method designed to redistribute
gradient mass away from low-dispersion groups and toward high-dispersion ones
produces exactly the predicted asymmetry across difficulty levels. The intended
operating regime---multi-task training where some tasks (e.g.\ rubric-graded
creative writing) sit permanently in the low-$\sigma_G$ regime that this paper
is designed to handle---is where the mechanism's value should be largest, and is
left to future work.

\section{Discussion}

\paragraph{Relation to standard GRPO.}
The proposal is a one-symbol change in Equation~\eqref{eq:grpo}: $\sigma_G
\rightarrow \hat{\sigma}_{\text{global}}(T_i)$, plus the clip in
Equation~\eqref{eq:clip}. No critic, no extra forward pass, no gradient surgery.
The EMA update in~\eqref{eq:ema} is $O(1)$ per batch per task.

\paragraph{Relation to prioritized sampling.}
The proportionality in Equation~\eqref{eq:budget} can be read as an implicit
priority weight: groups with larger-than-typical internal dispersion are
upweighted, smaller-than-typical groups are downweighted. Unlike prioritized
experience replay, the weighting is a free byproduct of changing the
normalization scale; no auxiliary buffer or sampler is required.

\paragraph{What is given up.}
The within-group advantages no longer have unit variance, so absolute advantage
magnitudes are no longer directly comparable to PPO conventions. Hyperparameters
that depend on advantage scale (notably the policy ratio clip and the KL
coefficient) may need light retuning, though in our setting the defaults from
GRPO transferred without modification. Tasks with persistent ceiling effects
(near-saturated benchmarks) will see slightly reduced learning---this is a
feature, not a bug, but practitioners targeting saturated-task accuracy should
be aware.

\paragraph{Limitations.}
The method assumes meaningful task partitions exist; for fully open-ended
training distributions a continuous task embedding would be needed in place of
discrete $T$. The EMA coefficient $\alpha$ trades responsiveness against
stability and is the one hyperparameter that genuinely matters; $\alpha=0.1$
worked across the settings we tried but a principled schedule is open work.
The MATH evaluation reported here is single-seed and does not control for
evaluation variance; multi-seed validation and the matching writing-task
evaluation are the natural next steps.
% =============================================================================
% V-GRPO: Theoretical Analysis Section
% Self-contained — drop into the v2 paper as a new Section 4 (or Appendix A).
% Required packages (already in v2): amsmath, amssymb, amsthm.
% =============================================================================

\section{Theoretical Analysis: V-GRPO as Implicit Entropy-Variance Regularization}
\label{sec:theory}

The empirical observation that V-GRPO maintains substantially higher policy
entropy than standard GRPO (Section~\ref{sec:results}) suggests that the method
acts, in some precise sense, as a regularizer. In this section we make this
intuition formal. We show that, under two interpretable assumptions about
reward--likelihood alignment and dispersion--entropy correlation, the V-GRPO
gradient is equivalent to the GRPO gradient minus the gradient of a specific
regularizer: the cross-prompt \emph{variance} of conditional policy entropy.
This is a strictly stronger claim than ``V-GRPO maximizes entropy,'' and
matches the empirical signature observed in our experiments more precisely
than a simple entropy-bonus framing.

\subsection{Setup and Notation}
\label{sec:theory:setup}

Fix a task type $T$ and a parameter setting $\theta$. Prompts are drawn from a
distribution $q \sim \mathcal{D}_T$. For each $q$, the current policy
$\pi_\theta(\cdot \mid q)$ is sampled $n$ times to produce a response group
$o_1, \dots, o_n$. Each response receives a reward $R_i = R(o_i \mid q)$. The
group statistics are
\begin{equation}
\mu_G(q) = \tfrac{1}{n}\sum_{i=1}^n R_i,
\qquad
\sigma_G^2(q) = \tfrac{1}{n}\sum_{i=1}^n (R_i - \mu_G(q))^2.
\end{equation}
Let $\hat\sigma^2(T) := \mathbb{E}_q[\sigma_G^2(q)]$ denote the converged
EMA-tracked global variance for task $T$, and define the relative dispersion
\begin{equation}
\rho_q := \frac{\sigma_G(q)}{\hat\sigma(T)},
\qquad \bar\rho := \mathbb{E}_q[\rho_q].
\end{equation}
Let $\mathcal{H}_q(\theta) := -\mathbb{E}_{o\sim\pi_\theta}[\log\pi_\theta(o\mid q)]$
denote the conditional policy entropy at prompt $q$, and
$\bar{\mathcal{H}} := \mathbb{E}_q[\mathcal{H}_q(\theta)]$ its prompt-averaged
value. Throughout we suppress the implicit dependence of $\mathcal{H}_q$ on
$\theta$ when it is clear from context.

The standard and V-GRPO advantages are
\begin{equation}
A_i^{\text{GRPO}} = \frac{R_i - \mu_G}{\sigma_G + \epsilon},
\qquad
A_i^{\text{V}} = \frac{R_i - \mu_G}{\hat\sigma(T) + \epsilon} = \rho_q \cdot A_i^{\text{GRPO}}.
\end{equation}
Ignoring the KL and ratio-clip terms, which are common to both methods and
contribute identical gradients, the core policy gradients are
\begin{align}
g_{\text{GRPO}}(\theta) &= \mathbb{E}_q\!\left[\sum_{i=1}^n A_i^{\text{GRPO}}
   \nabla_\theta \log\pi_\theta(o_i \mid q)\right], \\
g_{\text{V}}(\theta) &= \mathbb{E}_q\!\left[\rho_q \sum_{i=1}^n A_i^{\text{GRPO}}
   \nabla_\theta \log\pi_\theta(o_i \mid q)\right].
\end{align}
We define the per-prompt gradient direction
\begin{equation}
h_q(\theta) := \sum_{i=1}^n A_i^{\text{GRPO}} \nabla_\theta \log\pi_\theta(o_i \mid q).
\end{equation}

\subsection{Two Interpretable Assumptions}
\label{sec:theory:assumptions}

Our analysis rests on two assumptions, each of which has both an empirical
interpretation and a falsifiable consequence.

\begin{assumption}[Reward--Likelihood Alignment]
\label{a1}
There exists $\beta > 0$ and a residual function $\epsilon_q : \mathcal{O} \to \mathbb{R}$
such that for all $q$ and all $o$ in the support of $\pi_\theta(\cdot\mid q)$,
\begin{equation}
A(o, q) = \beta \log\pi_\theta(o \mid q) + \epsilon_q(o),
\label{eq:a1}
\end{equation}
where $A(o, q) := (R(o, q) - \mathbb{E}_{o'}[R(o',q)])/\sigma_G(q)$ is the
population advantage. We further write
$r_q := \mathbb{E}_{o\sim\pi_\theta}[\epsilon_q(o) \nabla_\theta \log\pi_\theta(o \mid q)]$
for the residual gradient contribution.
\end{assumption}

\begin{remark}
Assumption~\ref{a1} is the formalization of the design objective of RL
training: high-reward responses become more likely under the policy. The
parameter $\beta$ measures how tightly aligned reward and log-likelihood are
in the current policy. After convergence, $\beta$ is positive; early in
training it can be small but is not negative on average. The residual term
$r_q$ captures reward-driven gradient contributions not explained by
likelihood alignment alone.
\end{remark}

\begin{assumption}[Dispersion--Entropy Correlation]
\label{a2}
There exists $\gamma > 0$ and a noise term $\zeta_q$ with $\mathbb{E}_q[\zeta_q] = 0$
and $\mathbb{E}_q[\zeta_q \cdot \mathcal{H}_q] = 0$ such that
\begin{equation}
\rho_q - \bar\rho = \gamma \cdot (\mathcal{H}_q - \bar{\mathcal{H}}) + \zeta_q.
\label{eq:a2}
\end{equation}
\end{assumption}

\begin{remark}
Assumption~\ref{a2} states that prompts on which the policy is more uncertain
(higher $\mathcal{H}_q$) tend to produce response groups with greater within-group
reward dispersion. This is intuitively justified because high policy entropy
yields diverse sampled responses, which tend to receive diverse rewards under
any reward model with non-trivial response-discrimination capacity. The
parameter $\gamma$ measures the strength of this correlation; it is small
when the reward model is insensitive to response variation (e.g., binary
reward on saturated tasks) and large when the reward model finely
discriminates among responses.
\end{remark}

\begin{assumption}[EMA Stationarity]
\label{a3}
The EMA tracker has converged, so that $\hat\sigma^2(T) = \mathbb{E}_q[\sigma_G^2(q)]$.
Equivalently, $\mathbb{E}_q[\rho_q^2] = 1$.
\end{assumption}

Assumption~\ref{a3} is satisfied to leading order after the calibration phase
described in Section~\ref{sec:method}, with deviations of order $\alpha$ per
update step.

\subsection{Main Result}
\label{sec:theory:main}

\begin{proposition}[V-GRPO as Implicit Entropy-Variance Regularizer]
\label{prop:main}
Under Assumptions~\ref{a1}, \ref{a2}, and \ref{a3}, the V-GRPO policy
gradient admits the decomposition
\begin{equation}
g_{\text{V}}(\theta)
= \bar\rho \cdot g_{\text{GRPO}}(\theta)
\;-\; \frac{n \beta \gamma}{2}\, \nabla_\theta \mathrm{Var}_q[\mathcal{H}_q(\theta)]
\;+\; \mathcal{R}(\theta),
\label{eq:main}
\end{equation}
where $\bar\rho \in (0, 1]$ is the mean relative dispersion,
$\mathrm{Var}_q[\mathcal{H}_q(\theta)] = \mathbb{E}_q[(\mathcal{H}_q - \bar{\mathcal{H}})^2]$
is the cross-prompt variance of conditional policy entropy, and
\begin{equation}
\mathcal{R}(\theta) = n \gamma \cdot \mathbb{E}_q\big[(\mathcal{H}_q - \bar{\mathcal{H}})\, r_q\big]
\label{eq:residual}
\end{equation}
is a reward-driven residual.
\end{proposition}

\begin{corollary}[Equivalent Regularized Objective]
\label{cor:objective}
Performing gradient ascent with $g_{\text{V}}$ is equivalent, to first order
in the regularization strength, to performing gradient ascent on the modified
objective
\begin{equation}
\mathcal{L}_{\text{V}}(\theta)
= \bar\rho \cdot \mathcal{L}_{\text{GRPO}}(\theta)
\;-\; \frac{n \beta \gamma}{2}\, \mathrm{Var}_q[\mathcal{H}_q(\theta)],
\label{eq:objective}
\end{equation}
modulo the path integral of $\mathcal{R}(\theta)$.
\end{corollary}

\begin{proof}[Proof of Proposition~\ref{prop:main}]
We proceed in three steps.

\medskip
\noindent\textbf{Step 1: Decomposition of $g_{\text{V}}$ into a
covariance.} Writing $\rho_q = \bar\rho + (\rho_q - \bar\rho)$,
\begin{align}
g_{\text{V}}(\theta)
&= \mathbb{E}_q[\rho_q\, h_q(\theta)] \nonumber \\
&= \bar\rho \cdot \mathbb{E}_q[h_q(\theta)] + \mathbb{E}_q[(\rho_q - \bar\rho)\, h_q(\theta)] \nonumber \\
&= \bar\rho \cdot g_{\text{GRPO}}(\theta) + \mathrm{Cov}_q\big(\rho_q,\, h_q(\theta)\big),
\label{eq:proof:cov}
\end{align}
where the covariance is taken componentwise.

\medskip
\noindent\textbf{Step 2: Identification of $\mathbb{E}_o[h_q]$ with the entropy gradient.}
Taking the expectation of $h_q(\theta)$ over response sampling:
\begin{equation}
\mathbb{E}_{o_{1:n}\sim\pi_\theta}[h_q(\theta)] = n \cdot \mathbb{E}_{o\sim\pi_\theta}[A(o,q)\, \nabla_\theta \log\pi_\theta(o\mid q)].
\end{equation}
Substituting Assumption~\ref{a1}:
\begin{align}
\mathbb{E}_o[h_q(\theta)]
&= n\beta\, \mathbb{E}_o[\log\pi_\theta(o\mid q)\, \nabla_\theta \log\pi_\theta(o\mid q)] + n\, r_q.
\label{eq:proof:hq}
\end{align}
By the score-function identity, for any $\theta$-dependent function $f$,
\begin{equation}
\nabla_\theta \mathbb{E}_o[f(o,q)] = \mathbb{E}_o[f(o,q)\, \nabla_\theta \log\pi_\theta(o\mid q)] + \mathbb{E}_o[\nabla_\theta f(o,q)].
\end{equation}
Setting $f(o,q) = \log\pi_\theta(o\mid q)$ yields
\begin{align}
\nabla_\theta \mathbb{E}_o[\log\pi_\theta(o\mid q)]
&= \mathbb{E}_o[\log\pi_\theta\, \nabla_\theta \log\pi_\theta] + \underbrace{\mathbb{E}_o[\nabla_\theta \log\pi_\theta]}_{=0},
\end{align}
since $\mathbb{E}_o[\nabla_\theta \log\pi_\theta(o\mid q)] = \nabla_\theta \mathbb{E}_o[1] = 0$.
Therefore
\begin{equation}
\mathbb{E}_o[\log\pi_\theta(o\mid q)\, \nabla_\theta \log\pi_\theta(o\mid q)] = -\nabla_\theta \mathcal{H}_q(\theta).
\label{eq:proof:score}
\end{equation}
Substituting (\ref{eq:proof:score}) into (\ref{eq:proof:hq}):
\begin{equation}
\mathbb{E}_o[h_q(\theta)] = -n\beta\, \nabla_\theta \mathcal{H}_q(\theta) + n\, r_q.
\label{eq:proof:hq-final}
\end{equation}

\medskip
\noindent\textbf{Step 3: Combining the covariance with the entropy decomposition.}
Substituting (\ref{eq:proof:hq-final}) and Assumption~\ref{a2} into the
covariance term of (\ref{eq:proof:cov}):
\begin{align}
\mathrm{Cov}_q(\rho_q, h_q)
&= \mathbb{E}_q\big[(\rho_q - \bar\rho)\, h_q(\theta)\big] \nonumber\\
&= \mathbb{E}_q\!\left[\big(\gamma(\mathcal{H}_q - \bar{\mathcal{H}}) + \zeta_q\big)
   \cdot \big(-n\beta \nabla_\theta \mathcal{H}_q + n r_q\big)\right] \nonumber\\
&= -n\beta\gamma\, \mathbb{E}_q[(\mathcal{H}_q - \bar{\mathcal{H}})\, \nabla_\theta \mathcal{H}_q] \nonumber\\
&\quad + n\gamma\, \mathbb{E}_q[(\mathcal{H}_q - \bar{\mathcal{H}})\, r_q]
   - n\beta\, \mathbb{E}_q[\zeta_q \nabla_\theta \mathcal{H}_q]
   + n\, \mathbb{E}_q[\zeta_q r_q].
\label{eq:proof:expand}
\end{align}
By Assumption~\ref{a2}, $\mathbb{E}_q[\zeta_q \mathcal{H}_q] = 0$, which we extend
to $\mathbb{E}_q[\zeta_q \nabla_\theta \mathcal{H}_q] = 0$ under the regularity
assumption that $\zeta_q$ is uncorrelated with $\theta$-derivatives of
$\mathcal{H}_q$ (equivalent to $\zeta_q$ depending on $q$ only through quantities
independent of $\theta$). The fourth term in (\ref{eq:proof:expand}) is
absorbed into $\mathcal{R}(\theta)$.

For the first term, observe that
\begin{align}
\mathbb{E}_q[(\mathcal{H}_q - \bar{\mathcal{H}})\, \nabla_\theta \mathcal{H}_q]
&= \mathbb{E}_q[\mathcal{H}_q \nabla_\theta \mathcal{H}_q] - \bar{\mathcal{H}}\, \nabla_\theta \bar{\mathcal{H}} \nonumber\\
&= \tfrac{1}{2}\nabla_\theta \mathbb{E}_q[\mathcal{H}_q^2] - \tfrac{1}{2}\nabla_\theta \bar{\mathcal{H}}^2 \nonumber\\
&= \tfrac{1}{2} \nabla_\theta\!\left(\mathbb{E}_q[\mathcal{H}_q^2] - \bar{\mathcal{H}}^2\right) \nonumber\\
&= \tfrac{1}{2} \nabla_\theta\, \mathrm{Var}_q[\mathcal{H}_q(\theta)].
\label{eq:proof:var-identity}
\end{align}
Substituting (\ref{eq:proof:var-identity}) into (\ref{eq:proof:expand}) and
collecting terms:
\begin{equation}
\mathrm{Cov}_q(\rho_q, h_q) = -\frac{n\beta\gamma}{2}\, \nabla_\theta\, \mathrm{Var}_q[\mathcal{H}_q(\theta)]
+ \mathcal{R}(\theta),
\end{equation}
with $\mathcal{R}(\theta)$ as defined in (\ref{eq:residual}). Combining with
(\ref{eq:proof:cov}) yields (\ref{eq:main}). \qed
\end{proof}

\subsection{Interpretation}
\label{sec:theory:interpretation}

Proposition~\ref{prop:main} identifies the precise sense in which V-GRPO
regularizes the GRPO objective. The regularizer is \emph{not} the policy
entropy itself, but the cross-prompt variance of conditional policy entropy.
The two are conceptually distinct in an important way:

\begin{itemize}\itemsep4pt
\item \textbf{Standard entropy regularization} ($-\lambda \mathbb{E}_q[\mathcal{H}_q]$)
penalizes the policy whenever its average entropy is low. It uniformly
encourages the policy to be less confident, regardless of which prompts are
involved.
\item \textbf{V-GRPO's implicit regularizer}
($-\frac{n\beta\gamma}{2}\mathrm{Var}_q[\mathcal{H}_q]$) penalizes
\emph{disagreement} in confidence across prompts. It specifically resists
mode collapse on individual prompts---in which $\mathcal{H}_q$ becomes much
smaller than $\bar{\mathcal{H}}$---while having no first-order effect on the
average level of confidence.
\end{itemize}

This distinction predicts a specific empirical signature, summarized in
Table~\ref{tab:theory-predictions}, which is consistent with our observed
results.

\begin{table}[h]
\centering
\small
\begin{tabular}{p{0.42\textwidth} p{0.50\textwidth}}
\toprule
\textbf{Predicted by Prop.~\ref{prop:main}} & \textbf{Empirical observation} \\
\midrule
Average policy entropy increases modestly, not dramatically, since the
regularizer targets variance rather than mean entropy.
& Final entropy: $0.0557 \to 0.1001$ ($\sim$$2\times$, but in absolute terms small). \\
\addlinespace
The most pronounced effect is on the lower tail of $\mathcal{H}_q$: prompts
that under GRPO would collapse to near-zero conditional entropy are
prevented from doing so.
& Fraction of batches with $\sigma_G = 0$: $\sim 0.70 \to 0.00$. These are
exactly the batches where mode collapse on individual prompts has occurred. \\
\addlinespace
Easy/saturated prompts lose accuracy because they are no longer permitted
to peak as sharply.
& MATH Level 1 accuracy: $94.28\% \to 91.30\%$. \\
\addlinespace
Hard prompts gain because their gradient signal is no longer compressed
relative to easy ones; the variance regularizer indirectly elevates the
weight of high-$\mathcal{H}_q$ prompts in the update.
& MATH Level 5 accuracy: $47.21\% \to 48.26\%$. \\
\bottomrule
\end{tabular}
\caption{Predictions of Proposition~\ref{prop:main} and corresponding
empirical observations from Sections~\ref{sec:results} and onward.}
\label{tab:theory-predictions}
\end{table}

\subsection{Discussion of Assumptions}
\label{sec:theory:discussion}

\paragraph{Validity of (A1).}
Assumption~\ref{a1} is a linearization of the reward--likelihood relationship.
A constant $\beta$ across prompts is an idealization; in practice $\beta$
varies with $q$ and with training step. The relevant requirement for
Proposition~\ref{prop:main} is that $\mathbb{E}_q[\beta_q]$ exists and is positive,
which holds throughout normal RL training. The factor $\beta\gamma$ in
(\ref{eq:main}) should accordingly be read as a sign-preserving aggregate
strength of regularization rather than a literal product of constants.

\paragraph{Validity of (A2).}
Assumption~\ref{a2} requires that within-group reward dispersion correlate
positively with policy entropy. This is a property of the reward model rather
than of V-GRPO: it holds whenever the reward model can distinguish responses
that the policy treats as distinct. For binary or low-resolution reward
models on saturated tasks, $\gamma$ can be small, in which case
Proposition~\ref{prop:main} predicts that V-GRPO's effect on those tasks
will be correspondingly muted. This matches the small magnitude of effects
observed at MATH Levels 2--4 in our experiments.

\paragraph{The residual $\mathcal{R}(\theta)$.}
The residual term in (\ref{eq:residual}) captures reward-driven gradient
contributions that correlate with prompt-level entropy. In well-conditioned
training regimes, $r_q$ is uncorrelated with $\mathcal{H}_q$ on average, and
$\mathcal{R}(\theta) \approx 0$. The residual becomes non-negligible when
reward-model errors correlate with policy uncertainty---for instance, when
the reward model is consistently miscalibrated on high-entropy prompts.
This is itself a falsifiable prediction: V-GRPO's deviation from
the ideal regularization (\ref{eq:objective}) should grow precisely when
reward-model calibration is prompt-dependent. We leave empirical
investigation of this effect to future work.

\paragraph{Connection to clipping.}
The advantage clip $A_i^{\text{final}} = \mathrm{clip}(A_i^{\text{V}}, -C, C)$
introduced in Section~\ref{sec:method} does not appear in
Proposition~\ref{prop:main}; the analysis is unchanged when the clip is
inactive (the typical regime under well-behaved reward models). Under
reward-model malfunction, clipping bounds the influence of $\mathcal{R}(\theta)$
by capping the magnitude of $\epsilon_q$, providing a concrete robustness
guarantee complementary to the regularization analysis above.
\bibliographystyle{plainnat}
\begin{thebibliography}{1}
\bibitem[Shao et al., 2024]{shao2024deepseekmath}
Z.~Shao, P.~Wang, Q.~Zhu, R.~Xu, J.~Song, M.~Zhang, Y.~K.~Li, Y.~Wu, and D.~Guo.
\newblock DeepSeekMath: Pushing the limits of mathematical reasoning in open
  language models.
\newblock \emph{arXiv:2402.03300}, 2024.
\end{thebibliography}

\end{document}